% F035 指数幂的运算法则｜05 易错点：每点一个对照图
\begin{trapbox}
    \textbf{\textcolor{CTrap}{易错1｜“两幂相乘”和“幂的乘方”不是同一件事}}\par
    \textbf{容易错：} 把 $(2^3)^2$ 写成 $2^{3+2}$。\par
    \textbf{正确：} $(2^3)^2=2^{3\times2}=2^6$；而 $2^3\cdot2^2=2^{3+2}=2^5$。
    \begin{center}
    \begin{tikzpicture}[font=\small]
        \node[draw=orange!70,fill=orange!8,rounded corners=3pt,inner sep=6pt] at (0,0) {$2^3\cdot2^2=2^5$};
        \node[draw=green!55!black,fill=green!10,rounded corners=3pt,inner sep=6pt] at (5.1,0) {$(2^3)^2=2^6$};
        \node at (2.55,-0.65) {并排相乘：加 \qquad 再次乘方：乘};
    \end{tikzpicture}

    \end{center}

    \medskip
    \textbf{\textcolor{CTrap}{易错2｜只有“相乘”才能直接合并指数}}\par
    \textbf{容易错：} $2^3+2^4=2^{3+4}$。\par
    \textbf{正确：} $2^3+2^4=8+16=24$，而 $2^{3+4}=128$，两者不相等。
    \begin{center}
    \begin{tikzpicture}[font=\small]
        \node[draw=red!55,fill=red!5,rounded corners=3pt,inner sep=5pt] at (0,0) {$2^3+2^4=24$};
        \node[draw=green!55!black,fill=green!10,rounded corners=3pt,inner sep=5pt] at (4.8,0) {$2^3\cdot2^4=128$};
        \node at (2.4,-0.72) {中间是“$+$”还是“$\times$”？先看符号！};
    \end{tikzpicture}

    \end{center}

    \medskip
    \textbf{\textcolor{CTrap}{易错3｜积可以分别乘方，和不可以}}\par
    \textbf{容易错：} $(2+3)^2=2^2+3^2$。\par
    \textbf{正确：} $(2+3)^2=25\ne13=2^2+3^2$；只有\textbf{乘积}才有 $(2\cdot3)^2=2^2\cdot3^2=36$。
    \begin{center}
    \begin{tikzpicture}[font=\small]
        \node[draw=red!55,fill=red!5,rounded corners=3pt,inner sep=5pt] at (0,0) {$(2+3)^2\ne2^2+3^2$};
        \node[draw=green!55!black,fill=green!10,rounded corners=3pt,inner sep=5pt] at (5.9,0) {$(2\cdot3)^2=2^2\cdot3^2$};
        \node at (2.95,-0.75) {括号里是加法，不能照搬乘法的规律};
    \end{tikzpicture}

    \end{center}

    \medskip
    \textbf{\textcolor{CTrap}{易错4｜底数不同，不一定完全没办法}}\par
    \textbf{错误合并：} $2^2\cdot3^3\ne6^5$。\par
    \textbf{正确办法：} 若\textbf{指数相同}，则 $2^3\cdot3^3=(2\cdot3)^3=6^3$；若\textbf{能化成同底}，则 $4^2\cdot2^3=2^4\cdot2^3=2^7$。
    \begin{center}
    \begin{tikzpicture}[font=\small]
        \node[draw=blue!55,fill=blue!7,rounded corners=3pt,inner sep=5pt] at (0,0) {同底：$2^4\cdot2^3=2^7$};
        \node[draw=green!55!black,fill=green!10,rounded corners=3pt,inner sep=5pt] at (6.2,0) {同指数：$2^3\cdot3^3=6^3$};
        \node at (3.1,-0.75) {先看能不能“统一底数”或“统一指数”};
    \end{tikzpicture}

    \end{center}

    \medskip
    \textbf{\textcolor{CTrap}{易错5｜底数为负或为零，先判断有没有意义}}\par
    \textbf{错误做法：} 认为 $(-2)^{1/2}$ 和 $0^{-1}$ 都可以当作实数来算。\par
    \textbf{正确：} $(-2)^{1/2}=\sqrt{-2}$ 在实数范围无意义，$0^{-1}=1/0$ 无意义；但 $(-2)^3=-8$ 可以算。
    \begin{center}
    \begin{tikzpicture}[font=\small]
        \node[draw=red!55,fill=red!5,rounded corners=3pt,inner sep=5pt] at (0,0) {$(-2)^{1/2}$：不行};
        \node[draw=red!55,fill=red!5,rounded corners=3pt,inner sep=5pt] at (4.0,0) {$0^{-1}$：不行};
        \node[draw=green!55!black,fill=green!10,rounded corners=3pt,inner sep=5pt] at (7.9,0) {$(-2)^3=-8$：可以};
        \node at (3.9,-0.75) {实数指数优先检查：底数是否大于 $0$};
    \end{tikzpicture}

    \end{center}
\end{trapbox}
